\section*{Hoofdstuk \ref{ch:b1:stereochemistry-in-depth} --- Stereochemie verdiept}

\begin{solution}{exo:b1:stereochemistry-in-depth:1}
\omterm{def:b1:stereochemistry-in-depth:isomers}{Constitutie-isomeren}; \omterm{def:b1:stereochemistry-in-depth:enantiomers}{enantiomeren}; \omterm{def:b1:stereochemistry-in-depth:enantiomers}{diastereomeren} (\omterm{def:b1:stereochemistry-in-depth:isomers}{stereo-isomeren} die
geen spiegelbeelden zijn); twee \omterm{def:b1:stereochemistry-in-depth:isomers}{conformaties} van één verbinding;
\omterm{def:b1:stereochemistry-in-depth:enantiomers}{diastereomeren} (\cip{R,S} is de mesovorm).
\end{solution}

\begin{solution}{exo:b1:stereochemistry-in-depth:2}
\ce{-OH} $>$ \ce{-NH2} $>$ \ce{-CH3} $>$ \ce{-H} (O, N, C, H).
\ce{-COOH} (O,O,O) $>$ \ce{-CHO} (O,O,H) $>$ \ce{-CH2OH} (O,H,H) $>$
\ce{-CH(CH3)2} (C,C,H). \ce{-Br} $>$ \ce{-Cl} $>$ \ce{-CCl3} (Cl,Cl,Cl)
$>$ \ce{-CH2Cl} (Cl,H,H): het eerste atoom beslist vóór zijn buren.
\end{solution}

\begin{solution}{exo:b1:stereochemistry-in-depth:3}
2,3-Dichloorbutaan: twee gelijkwaardige centra, \cip{R,R}, \cip{S,S} en
het meso-\cip{R,S}: drie. 2-Broom-3-chloorbutaan: twee verschillende
centra, vier \omterm{def:b1:stereochemistry-in-depth:isomers}{stereo-isomeren}, geen mesovorm. Pentaan-2,4-diol: drie,
waarvan \cip{2R,4S} meso is.
\end{solution}

\begin{solution}{exo:b1:stereochemistry-in-depth:4}
Op C5: isopropenyl (C,C,C, waarbij de dubbele binding twee keer telt)
$>$ C6 $>$ C4 $>$ H. C6 en C4 dragen allebei (C,H,H); één stap verder
leidt C6 naar het carbonylkoolstofatoom (O,O,C) en C4 naar C3 (C,C,H): C6
wint. Met de isopropenylgroep vooraan en H achteraan draait de volgorde
isopropenyl $\to$ C6 $\to$ C4 zoals getekend met de wijzers van de klok
mee: \cip{R}, het carvon van groene munt; het gearceerde is \cip{S}.
\end{solution}

\begin{solution}{exo:b1:stereochemistry-in-depth:5}
$\theta = 0$: methylgroepen eclips, \qty{16.5}{kJ/mol}; $60^\circ$: gauche,
ongeveer \qty{2.9}{kJ/mol} (minimum 2.8 bij $62^\circ$); $120^\circ$:
methyl tegenover H, \qty{15.2}{kJ/mol}; $180^\circ$: anti, 0. Met gelijke
energieën zijn er twee \omterm{def:b1:stereochemistry-in-depth:conformations}{gaucheconformaties} voor één anti: twee derde
gauche.
\end{solution}

\begin{solution}{exo:b1:stereochemistry-in-depth:6}
Keten verticaal met \ce{COOH} bovenaan en \ce{CH3} onderaan; rangorde
$\ce{NH2} > \ce{COOH} > \ce{CH3} > \ce{H}$. Met \ce{NH2} links en H
rechts draait $\ce{NH2} \to \ce{COOH} \to \ce{CH3}$ zoals gezien met de
wijzers van de klok mee; H staat horizontaal (naar de lezer toe), dus de
descriptor wordt omgekeerd: \cip{S}. De aminogroep staat links.
\end{solution}

\begin{solution}{exo:b1:stereochemistry-in-depth:7}
\textit{cis}: op naburige koolstofatomen geeft één binding omhoog en één
omlaag onder de \omterm{def:b1:stereochemistry-in-depth:conformations}{axiale} en equatoriale posities in beide \omterm{def:b1:stereochemistry-in-depth:conformations}{stoelen}
axiaal--equatoriaal. \textit{trans}: \omterm{def:b1:stereochemistry-in-depth:conformations}{di-equatoriaal} of diaxiaal. In de
\omterm{def:b1:stereochemistry-in-depth:conformations}{di-equatoriale} \omterm{def:b1:stereochemistry-in-depth:conformations}{stoel} staan de twee methylgroepen gauche ten opzichte van
elkaar (één interactie); het \textit{cis}-isomeer heeft die interactie
plus een \omterm{def:b1:stereochemistry-in-depth:conformations}{axiale} methylgroep, ongeveer \qty{5.7}{kJ/mol} meer. Het
\textit{trans}-isomeer is het stabielst.
\end{solution}

\begin{solution}{exo:b1:stereochemistry-in-depth:8}
$[\alpha] = 13.2/(2.00 \times 0.100) = +66.0^\circ$. Een buis van
\qty{1.00}{dm}: $+6.60^\circ$. Het \omterm{def:b1:stereochemistry-in-depth:enantiomers}{enantiomeer}: $-13.2^\circ$ in de buis
van \qty{2.00}{dm}.
\end{solution}

\begin{solution}{exo:b1:stereochemistry-in-depth:9}
\qty{25}{\celsius}: $\exp(5700/(8.314 \times 298.15)) = 10.0$, fractie
$10.0/11.0 = \qty{91}{\percent}$. \qty{-50}{\celsius}: $\exp(5700/(8.314
\times 223.15)) = 21.6$, \qty{96}{\percent}.
\end{solution}

\begin{solution}{exo:b1:stereochemistry-in-depth:10}
In ethaan zijn alle zes waterstofatomen gelijk: elke \omterm{def:b1:stereochemistry-in-depth:conformations}{gestaffelde
conformatie} is dezelfde. In butaan kan de achterste methylgroep tegenover
de voorste staan (anti) of ernaast ($\pm 60^\circ$, gauche), waar de twee
methylgroepen elkaar in de weg zitten. Populaties: anti 1, gauche
$2\exp(-2.8/2.48) = 0.65$: ongeveer \qty{61}{\percent} anti en
\qty{39}{\percent} gauche.
\end{solution}

\begin{solution}{exo:b1:stereochemistry-in-depth:11}
C2 en C4 zijn stereogeen. C3 draagt H, OH en twee takken die qua
constitutie identiek zijn; ze verschillen alleen wanneer C2 en C4
tegengestelde descriptoren hebben, en alleen dan is C3 stereogeen
(pseudo-asymmetrisch). Isomeren: \cip{2R,4R} en \cip{2S,4S}, een paar
\omterm{def:b1:stereochemistry-in-depth:enantiomers}{enantiomeren} (C3 niet stereogeen); \cip{2R,4S} met de twee schikkingen
bij C3, twee mesovormen. Vier in totaal.
\end{solution}

\begin{solution}{exo:b1:stereochemistry-in-depth:12}
$[\alpha]_{\text{monster}} = 2.31/(1.00 \times 0.050) = +46.2^\circ$, en
$2x - 1 = 46.2/66.0 = 0.700$: $x = 0.850$, \qty{85}{\percent} van het
$(+)$-enantiomeer en \qty{15}{\percent} van het $(-)$-enantiomeer.
\end{solution}

\begin{solution}{pb:b1:stereochemistry-in-depth:1}
\textbf{1.} Cyclohexaanring: C1 draagt OH, C2 de isopropylgroep, C5 de
methylgroep; C1, C2 en C5 zijn stereogeen.
\textbf{2.} $2^3 = 8$, vier paren \omterm{def:b1:stereochemistry-in-depth:enantiomers}{enantiomeren}.
\textbf{3.} Geen enkele schikking heeft een spiegelvlak: de drie
substituenten zijn allemaal verschillend.
\textbf{4.} \ce{OH} $>$ C2 (C,C,H) $>$ C6 (C,H,H) $>$ H.
\textbf{5.} C1 (O,C,H) $>$ isopropyl (C,C,H) $>$ C3 (C,H,H) $>$ H.
\textbf{6.} C6 en C4 allebei (C,H,H); daarna leidt C6 naar C1 (O,C,H), C4
naar C3 (C,H,H): C6 $>$ C4 $>$ \ce{CH3} $>$ H.
\textbf{7.} \cip{1S,2R,5S}.
\textbf{8.} In een \omterm{def:b1:stereochemistry-in-depth:conformations}{stoel} wijzen twee \omterm{def:b1:stereochemistry-in-depth:conformations}{equatoriale} bindingen op naburige
koolstofatomen (C1, C2) de ene omhoog en de andere omlaag:
\textit{trans}; op de koolstofatomen 1 en 3 van de ring (C1 en C5, via
C6) wijzen beide \omterm{def:b1:stereochemistry-in-depth:conformations}{equatoriale} bindingen dezelfde kant op: \textit{cis}.
Dat is de schikking van menthol.
\textbf{9.} Alle drie worden \omterm{def:b1:stereochemistry-in-depth:conformations}{axiaal}.
\textbf{10.} In de omgeklapte \omterm{def:b1:stereochemistry-in-depth:conformations}{stoel} zou de methylgroep alleen al ongeveer
\qty{5.7}{kJ/mol} kosten en de grotere isopropylgroep meer; de OH voegt
een kleinere eigen kostprijs toe, en de \omterm{def:b1:stereochemistry-in-depth:conformations}{axiale} OH en methylgroep, op de
koolstofatomen 1 en 3 en aan dezelfde kant, zouden elkaar ook in de weg
zitten. Dat is ruim meer dan \qty{12}{kJ/mol}, een verhouding van meer
dan $\exp(12000/2479) \approx 130$ op één: menthol zit bijna volledig in
de \omterm{def:b1:stereochemistry-in-depth:conformations}{stoel} met alles \omterm{def:b1:stereochemistry-in-depth:conformations}{equatoriaal}.
\textbf{11.} Een \omterm{def:b1:stereochemistry-in-depth:enantiomers}{diastereomeer}: maar één van de drie centra verschilt.
\textbf{12.} Met isopropyl en methyl \omterm{def:b1:stereochemistry-in-depth:conformations}{equatoriaal} staat de OH \omterm{def:b1:stereochemistry-in-depth:conformations}{axiaal}.
\textbf{13.} \omterm{def:b1:stereochemistry-in-depth:enantiomers}{Diastereomeren} verschillen in al hun eigenschappen (hier de
\omterm{def:b1:intermolecular-forces-solvents:hbond}{waterstofbruggen} van een \omterm{def:b1:stereochemistry-in-depth:conformations}{axiale} of \omterm{def:b1:stereochemistry-in-depth:conformations}{equatoriale} OH); \omterm{def:b1:stereochemistry-in-depth:enantiomers}{enantiomeren}
verschillen alleen tegenover \omterm{def:b1:stereochemistry-in-depth:stereocentre}{chirale} partners.
\textbf{14.} $[\alpha] = -2.46/(1.00 \times 0.0500) = -49.2^\circ$, binnen
het bereik van het handboek.
% ledger: alpha:l-menthol-range
\textbf{15.} $+49.2^\circ$; $0$.
\textbf{16.} $-1.97/0.0500 = -39.4^\circ$.
\textbf{17.} $\alpha = \ell c[\alpha]_{(-)}\bigl(x - (1 - x)\bigr) =
\ell c[\alpha]_{(-)}(2x - 1)$.
\textbf{18.} $x = \frac12(1 + \alpha/\alpha_{\text{ref}})$ bij gelijke $\ell$ en $c$.
\textbf{19.} Ja: elke andere optisch actieve verbinding voegt haar eigen
term toe (\omterm{prop:b1:stereochemistry-in-depth:biot}{wet van Biot}). Een controle van de chemische zuiverheid met een
andere methode (chromatografie) is nodig voordat de rotatie als een
samenstelling gelezen wordt.
\textbf{20.} $x = \frac12(1 + 1.97/2.46) = 0.900$.
\textbf{21.} $x = \frac12(1 + 2.40/2.46) = 0.988$.
\textbf{22.} Het tweede (\qty{98.8}{\percent}).
\textbf{23.} $-3.94^\circ$; dezelfde afleesfout is dan een kleinere
fractie van de hoek.
\textbf{24.} \textbf{\qty{90.0}{\percent} $(-)$-menthol} (en
\qty{10.0}{\percent} van zijn \omterm{def:b1:stereochemistry-in-depth:enantiomers}{enantiomeer}).
\end{solution}
